\documentclass[11pt]{article}
\usepackage[margin=1in]{geometry}
\usepackage{amsmath,amssymb,amsthm}
\usepackage{graphicx}
\usepackage[hidelinks]{hyperref}
\setlength{\emergencystretch}{2em}

\newtheorem{theorem}{Theorem}
\newtheorem{lemma}{Lemma}
\newtheorem{corollary}{Corollary}

\newcommand{\C}{\mathbb C}

\title{Product Quotient \(C^*\)-Algebras Fail the Ball-Covering Property}
\author{Run \texttt{fa\_banach\_001}, lane 19}
\date{June 17, 2026}

\begin{document}
\maketitle

\section*{Source Problem}

The source paper is Jinghao Huang, Karimbergen Kudaybergenov, and Rui Liu,
\emph{The ball-covering property of von Neumann algebras and noncommutative
symmetric spaces}, arXiv:2311.14257.  After proving that a commutative unital
\(C^*\)-algebra \(C(K)\) has BCP/UBCP exactly when its pure state space, i.e.
\(K\), has a countable \(\pi\)-basis, the authors ask for a noncommutative
version:

\begin{center}
\includegraphics[width=0.92\textwidth]{figures/open_problem_crop.png}
\end{center}

\section*{Statement}

Let \((E_n)\) be a sequence of nonzero normed spaces.  Write
\[
        \ell_\infty(E_n)=\{x=(x_n):\sup_n\|x_n\|<\infty\}
\]
and
\[
        c_0(E_n)=\{x=(x_n):\|x_n\|\to0\}.
\]
For \(x\in\ell_\infty(E_n)\), denote its quotient class by \([x]\).  The
quotient norm is
\[
        \|[x]\|=\limsup_{n\to\infty}\|x_n\|.
\]

\begin{theorem}
\label{thm:banach-product}
For every sequence \((E_n)\) of nonzero normed spaces,
\[
        \ell_\infty(E_n)/c_0(E_n)
\]
fails BCP.
\end{theorem}

\begin{corollary}
\label{cor:cstar-product}
Let \((A_n)\) be a sequence of nonzero \(C^*\)-algebras.  Then the
\(C^*\)-algebra
\[
        \prod_{n=1}^\infty A_n\Big/\bigoplus_{n=1}^\infty A_n
\]
fails BCP.  In particular, for every nonzero \(C^*\)-algebra \(A\), the
sequence-algebra quotient \(\ell_\infty(A)/c_0(A)\) fails BCP.
\end{corollary}

\begin{corollary}[quotient instability]
\label{cor:quotient-instability}
BCP and UBCP are not preserved by \(C^*\)-quotients.  More precisely, if
\(A\) is any nonzero separable \(C^*\)-algebra, then \(\ell_\infty(A)\) has
UBCP, but its quotient \(\ell_\infty(A)/c_0(A)\) fails BCP.
\end{corollary}

\section*{Idea of the Proof}

This is the diagonal escape proof for \(\ell_\infty/c_0\), but with arbitrary
nonzero fibers.  Given countably many proposed centers in the quotient, choose
for each center infinitely many coordinates where it nearly attains its
quotient norm.  On the selected coordinates, choose a unit vector in the fiber
that is far enough from the corresponding center coordinate.  The resulting
bounded sequence has quotient norm one, while every proposed center sees
infinitely many coordinates on which the distance is at least its own norm.

\section*{Proof}

We use the standard equivalent form of BCP: a normed space \(Y\) has BCP iff
there are countably many centers \(y_k\in Y\) such that
\[
        S_Y\subseteq \bigcup_{k=1}^\infty B(y_k,\|y_k\|).
\]

\begin{lemma}[one-coordinate escape]
\label{lem:one-coordinate}
Let \(E\ne\{0\}\) be a normed space and let \(y\in E\).  For every
\(\varepsilon>0\), there is \(u\in S_E\) such that
\[
        \|u-y\|\ge \|y\|+1-\varepsilon.
\]
In particular \(\|u-y\|\ge \|y\|\) for a suitable choice of \(u\).
\end{lemma}

\begin{proof}
If \(y=0\), take any \(u\in S_E\).  If \(y\ne0\), choose \(f\in S_{E^*}\) with
\(f(y)=\|y\|\).  Since \(\|f\|=1\), after changing phase if necessary one can
pick \(u\in S_E\) with \(\operatorname{Re} f(u)<-1+\varepsilon\).
Then
\[
        \|u-y\|\ge \operatorname{Re} f(y-u)>\|y\|+1-\varepsilon.
\]
Taking, say, \(\varepsilon=1/2\) gives the final assertion.
\end{proof}

\begin{lemma}
\label{lem:repeat-map}
There is a map \(\pi:\mathbb N\to\mathbb N\) such that for every \(i\) and
every \(N\), some \(k\ge N\) satisfies \(\pi(k)=i\).
\end{lemma}

\begin{proof}
List \(\mathbb N\times\mathbb N\) as a sequence \((i_k,j_k)\), and put
\(\pi(k)=i_k\), arranging the list so each first coordinate occurs infinitely
often.
\end{proof}

\begin{proof}[Proof of Theorem~\ref{thm:banach-product}]
Let \(\xi_i=[a^i]\), \(i\ge1\), be any countable family in
\(\ell_\infty(E_n)/c_0(E_n)\), with representatives
\[
        a^i=(a^i_n)_{n\ge1}.
\]
We construct \(b=(b_n)\in\ell_\infty(E_n)\) such that \(\|[b]\|=1\) and
\[
        \|[b]-\xi_i\|\ge \|\xi_i\|,\qquad i\ge1.
\]

For each \(i\), choose an increasing sequence \((n_{i,r})_{r\ge1}\) such that
\[
        \|a^i_{n_{i,r}}\|\longrightarrow \|\xi_i\|.
\]
Let \(\pi\) be as in Lemma~\ref{lem:repeat-map}.  We recursively choose a
strictly increasing sequence of coordinates \(m_k\).  Having chosen
\(m_{k-1}\), choose \(m_k=n_{\pi(k),r}\) for some large \(r\) so that
\[
        m_k>m_{k-1}
        \quad\text{and}\quad
        \|a^{\pi(k)}_{m_k}\|>\|\xi_{\pi(k)}\|-1/k
\]
whenever \(\|\xi_{\pi(k)}\|>0\).  If \(\|\xi_{\pi(k)}\|=0\), only require
\(m_k>m_{k-1}\).  In all cases \(m_k\) is chosen from the fixed subsequence
\((n_{\pi(k),r})_r\).

At coordinate \(m_k\), apply Lemma~\ref{lem:one-coordinate} in the nonzero
space \(E_{m_k}\) to choose \(b_{m_k}\in S_{E_{m_k}}\) such that
\[
        \|b_{m_k}-a^{\pi(k)}_{m_k}\|\ge \|a^{\pi(k)}_{m_k}\|.
\]
For all other coordinates, set \(b_n=0\).  Then \(\|b\|_{\ell_\infty}=1\), and
because the selected coordinates are infinite, \(\|[b]\|=1\).

Fix \(i\).  Since \(\pi(k)=i\) for infinitely many \(k\), along that infinite
subsequence the selected coordinates \(m_k\) are drawn from \((n_{i,r})_r\)
with \(r\to\infty\), and hence
\[
        \|a^i_{m_k}\|\longrightarrow \|\xi_i\|.
\]
Together with
\[
        \|b_{m_k}-a^i_{m_k}\|\ge \|a^i_{m_k}\|
\]
this gives
\[
        \|[b]-\xi_i\|
        =
        \limsup_n\|b_n-a^i_n\|
        \ge \|\xi_i\|.
\]
Thus \([b]\in S_{\ell_\infty(E_n)/c_0(E_n)}\) is in none of the open balls
\(B(\xi_i,\|\xi_i\|)\).  Since the countable family \((\xi_i)\) was arbitrary,
the quotient fails BCP.
\end{proof}

Corollary~\ref{cor:cstar-product} follows by taking \(E_n=A_n\) with its \(C^*\)-norm.  The
quotient \(\prod A_n/\bigoplus A_n\) is exactly
\(\ell_\infty(A_n)/c_0(A_n)\) as a Banach space.

\begin{proof}[Proof of Corollary~\ref{cor:quotient-instability}]
It remains only to prove that \(\ell_\infty(A)\) has UBCP when \(A\) is
separable and nonzero.  Choose a countable dense set \(D\subseteq S_A\).  For
each \(n\in\mathbb N\) and \(d\in D\), define \(c^{n,d}\in \ell_\infty(A)\) by
\[
        c^{n,d}_j=
        \begin{cases}
        4d,& j=n,\\
        0,& j\ne n.
        \end{cases}
\]
Then \(\|c^{n,d}\|=4\), so every ball \(B(c^{n,d},7/2)\) avoids the origin.
Let \(x=(x_j)\in S_{\ell_\infty(A)}\).  Pick \(n\) with
\(\|x_n\|>3/4\), put \(\alpha=\|x_n\|\), and set \(u=x_n/\alpha\).  Choose
\(d\in D\) with \(\|d-u\|<1/64\).  At the \(n\)-th coordinate,
\[
        \|x_n-4d\|
        \le 4-\alpha+4\|u-d\|
        < 4-\frac34+\frac1{16}
        < \frac72.
\]
At all other coordinates, \(\|x_j-c^{n,d}_j\|=\|x_j\|\le1<7/2\).  Hence
\(\|x-c^{n,d}\|<7/2\), and the countable family of equal-radius balls
\(B(c^{n,d},7/2)\) covers \(S_{\ell_\infty(A)}\).  Thus \(\ell_\infty(A)\)
has UBCP, while Corollary~\ref{cor:cstar-product} says
\(\ell_\infty(A)/c_0(A)\) fails BCP.
\end{proof}

\section*{Scope and Relation to the C*-Algebra Question}

This gives a broad negative family for the source question.  It does not
characterize arbitrary \(C^*\)-algebras, but it shows that sequence-algebra
coronas are robustly outside BCP.  In particular, even when the fibers are
separable or finite dimensional, the product quotient has too much tail room
for any countable ball cover to capture its unit sphere.  The quotient
instability corollary is also a warning: a C*-algebraic characterization of
BCP cannot be based on quotient-stable invariants alone.

\section*{Novelty and Search Bounds}

The run indexes were searched for arXiv:2311.14257, product quotients,
\(\ell_\infty/c_0\), Calkin/corona algebras, and BCP.  Web and local source
checks found the classical scalar obstruction and the later Calkin-algebra
paper \cite{BLS}, but no exact product-of-arbitrary-fibers statement in the
run memory.  The theorem is a direct generalization of the scalar
\(\ell_\infty/c_0\) diagonal construction.

\section*{Human Review Focus}

The proof is elementary.  Review the one-coordinate escape lemma and the
diagonal construction of the coordinates \(m_k\), especially the fact that
each center \(i\) is attacked along infinitely many coordinates where its
representative nearly realizes the quotient norm.

\begin{thebibliography}{9}

\bibitem{HKL}
Jinghao Huang, Karimbergen Kudaybergenov, and Rui Liu,
\emph{The ball-covering property of von Neumann algebras and noncommutative
symmetric spaces},
arXiv:2311.14257, 2023.

\bibitem{BLS}
Qiyao Bao, Rui Liu, and Jie Shen,
\emph{The ball-covering property of non-commutative spaces of operators on
Banach spaces},
arXiv:2412.07137, 2024.

\end{thebibliography}

\end{document}
